\documentclass[11pt]{article}
\usepackage[margin=2.5cm]{geometry}
\usepackage{amsmath}
\usepackage[colorlinks,citecolor=blue,urlcolor=blue]{hyperref}
\usepackage{booktabs}

\title{Simulations of Phytic Acid near Oleosin Covered Oil Drop\\
       \large Working document}
\author{}
\date{}

\begin{document}
\maketitle

\section*{Methods draft --- \S2.2.10 Computational simulations}

Coarse-grained Monte Carlo simulations resolved the phytate--oleosin
interaction at the molecular scale.  The N-terminal arm of
\textit{Cannabis sativa} oleosin (UniProt~A0A803P977; sequence
\texttt{MKSFLPESGPSAQQ}) was represented at one bead per residue with the
CALVADOS~3 model~\cite{Tesei2021,Cao2024}.  Fifty such chains were end-grafted to
one wall of a rectangular slit ($130\times130\times200$\,\AA$^{3}$,
periodic in $x$ and $y$) at a surface density of
$3\times10^{-3}$\,\AA$^{-2}$, forming a brush that represents the
oleosin-covered oil-body surface~\cite{Tzen1992}.  A single phytate
molecule (\textit{myo}-inositol hexakisphosphate) was modelled as a rigid
seven-bead unit: one inositol centre and six phosphate groups in the chair
geometry.  Beads interacted through the Ashbaugh--Hatch
potential~\cite{Ashbaugh2008} and a screened Coulomb (Debye--H\"uckel) term
set by the salt concentration, with phytate--peptide cross-interactions
made purely repulsive to suppress spurious hydrophobic attraction.  The
solvent entered implicitly through the dielectric permittivity of water at
298.15\,K.

Peptide and phytate both titrated with pH.  The ionisable peptide groups
(Glu, Lys, and the two chain termini) and all twelve phytate protons
exchanged with an implicit proton bath ($a_{\mathrm{H}^+}=10^{-\mathrm{pH}}$)
through atom-type and molecular swap moves in the semi-grand canonical
ensemble~\cite{SmithTriska1994}, with p$K_\mathrm{a}$ values taken from
potentiometric titrations~\cite{Grimsley2009,DeStefano2003}.  Configurations
were sampled by Metropolis Monte Carlo in Faunus~\cite{Stenqvist2013} using
translation, rotation, and pivot moves, at pH~4.0, 6.0, 7.0, and 10.0 and
ionic strengths of 5--150\,mM NaCl.

The free energy of phytate as a function of distance from the brush was
obtained by umbrella sampling along the inositol $z$-coordinate, using 40
overlapping windows that spanned the grafted wall to the slit centre.  A
harmonic drive first steered the molecule into each window, the bias was
then removed for production sampling, and neighbouring windows were combined
through their histogram overlap.  This potential of mean force is reported
against $r$, the distance from the grafted wall, and referenced to zero in
the bulk.  The same simulations also yielded the mean phytate charge, its
rotational mean-square displacement (a relative measure of rotational
mobility), and the laterally averaged electric potential $\varphi(r)$, from
which a $\zeta$-potential was estimated as $\varphi$ at the outer edge of the
grafted charge.  Finally, the rotational and translational diffusion
coefficients of free phytate were estimated for each protonation state from
RDKit conformer ensembles~\cite{RDKit} using Rotne--Prager--Yamakawa bead
hydrodynamics~\cite{Zuk2018,Waszkiewicz2023}.

\bigskip
\hrule
\bigskip
\noindent\textit{Detailed working notes follow.}

\section{System}

We model the N-terminal arm of \textit{Cannabis sativa} oleosin
(UniProt A0A803P977) grafted to a planar surface that mimics the oil
body monolayer~\cite{Tzen1992}.  The FASTA sequence of the modelled
fragment is \texttt{MKSFLPESGPSAQQ}, where lowercase \texttt{n}
and \texttt{c} denote coarse-grained N- and C-terminal beads and
\texttt{a} is an anchor bead tethered to the surface.
The 14 amino acid residues are represented at one-bead-per-residue
resolution using the CALVADOS~3 force field~\cite{Tesei2021,Cao2024}.
Consecutive beads are connected by harmonic bonds with force constant
$k = 80.33$\,kJ\,mol$^{-1}$\,\AA$^{-2}$ and equilibrium distance
$r_\text{eq} = 3.8$\,\AA.

$N = 50$ chains are grafted at surface density
$\sigma \approx 0.003$\,\AA$^{-2}$ (brush regime) on the bottom wall
of a slit geometry ($130 \times 130 \times 200$\,\AA$^{3}$, periodic
in $xy$, hard walls at $z = \pm 100$\,\AA).  Grafting is realised by
a harmonic tether on the anchor bead at $z_0 \approx -100$\,\AA\ with
$k = 0.5$\,kJ\,mol$^{-1}$\,\AA$^{-2}$.

\section{Peptide titration}

Four residue types undergo pH-dependent protonation via atom-type swap
moves with implicit protons ($a_{\text{H}^+} = 10^{-\text{pH}}$).
The protonated and deprotonated forms differ only in charge:

\begin{table}[h]
\centering
\begin{tabular}{@{}lllc@{}}
\toprule
Residue & Protonated (charge) & Deprotonated (charge) & $\text{p}K_\text{a}$ \\
\midrule
GLU  & HGLU (0)   & GLU ($-1$)  & 4.24 \\
LYS  & LYS ($+1$) & LYS0 (0)   & 10.4 \\
NTR  & HNTR ($+1$) & NTR (0)   & 8.0  \\
CTR  & HCTR (0)   & CTR ($-1$)  & 3.67 \\
\bottomrule
\end{tabular}
\caption{Peptide residue titration reactions.
  $\text{p}K_\text{a}$ values from Ref.~\cite{Grimsley2009}.}
\end{table}

\noindent
At pH~7 the expected mean charge per chain is approximately $-0.1$\,e
(NTR mostly protonated, LYS protonated, GLU deprotonated, CTR deprotonated).
At pH~4 the brush becomes net positive ($\approx +1$\,e per chain);
at pH~10 it becomes strongly negative ($\approx -1.3$\,e per chain).

\section{Phytate titration}

One molecule of myo-inositol hexakisphosphate (phytate, IP6) is
included as a titrating ligand.  Each molecule is coarse-grained as
seven beads: one inositol centre (INO, $\sigma = 6.2$\,\AA,
$M = 180$\,Da) and six phosphate groups (PH, $\sigma = 5.8$\,\AA,
$M = 95$\,Da) placed at the crystallographic P positions in the
chair conformation ($\approx 4.5$\,\AA\ from the ring centre, one
axial and five equatorial).  Bead sizes are scaled from CALVADOS~3
by $M^{1/3}$ to maintain volume consistency with the amino acid beads.

Phytate has 12 dissociable protons (two per phosphate group),
giving 13 charge states from H$_{12}$PA (neutral) to PA$^{12-}$.
All 12 dissociation steps are modelled via molecular swap moves
in the semi-grand canonical ensemble~\cite{SmithTriska1994}.
Macroscopic $\text{p}K_\text{a}$ values are taken from potentiometric
measurements in NaCl solution~\cite{DeStefano2003}:

\begin{table}[h]
\centering
\begin{tabular}{@{}lcc@{}}
\toprule
Reaction & $\text{p}K_\text{a}$ & Charge change \\
\midrule
H$_{12}$PA       $\rightleftharpoons$ H$_{11}$PA$^{1-}$  + H$^+$ & 1.1   & $0 \to -1$   \\
H$_{11}$PA$^{1-}$ $\rightleftharpoons$ H$_{10}$PA$^{2-}$ + H$^+$ & 1.5   & $-1 \to -2$  \\
H$_{10}$PA$^{2-}$ $\rightleftharpoons$ H$_9$PA$^{3-}$    + H$^+$ & 1.7   & $-2 \to -3$  \\
H$_9$PA$^{3-}$   $\rightleftharpoons$ H$_8$PA$^{4-}$     + H$^+$ & 2.1   & $-3 \to -4$  \\
H$_8$PA$^{4-}$   $\rightleftharpoons$ H$_7$PA$^{5-}$     + H$^+$ & 2.5   & $-4 \to -5$  \\
H$_7$PA$^{5-}$   $\rightleftharpoons$ H$_6$PA$^{6-}$     + H$^+$ & 2.9   & $-5 \to -6$  \\
H$_6$PA$^{6-}$   $\rightleftharpoons$ H$_5$PA$^{7-}$     + H$^+$ & 5.72  & $-6 \to -7$  \\
H$_5$PA$^{7-}$   $\rightleftharpoons$ H$_4$PA$^{8-}$     + H$^+$ & 6.28  & $-7 \to -8$  \\
H$_4$PA$^{8-}$   $\rightleftharpoons$ H$_3$PA$^{9-}$     + H$^+$ & 6.81  & $-8 \to -9$  \\
H$_3$PA$^{9-}$   $\rightleftharpoons$ H$_2$PA$^{10-}$    + H$^+$ & 7.60  & $-9 \to -10$ \\
H$_2$PA$^{10-}$  $\rightleftharpoons$ HPA$^{11-}$        + H$^+$ & 9.94  & $-10 \to -11$ \\
HPA$^{11-}$      $\rightleftharpoons$ PA$^{12-}$         + H$^+$ & 11.84 & $-11 \to -12$ \\
\bottomrule
\end{tabular}
\end{table}

\noindent
The net charge is distributed uniformly over the six phosphate beads;
protons are treated implicitly with activity $a_{\text{H}^+} =
10^{-\text{pH}}$.  All 13 protonation states share the same molecular
geometry (seven beads), enabling molecular swap moves that preserve
spatial orientation via gyration-tensor alignment.

\section{Interactions}

Nonbonded interactions between all beads combine the Ashbaugh--Hatch
potential~\cite{Ashbaugh2008} (short-range, hydrophobicity-dependent)
with screened Coulomb (Debye--H\"uckel) electrostatics.
For phytate--amino acid cross-pairs, the Ashbaugh--Hatch
$\lambda$ parameter is set to zero (purely repulsive WCA) to avoid
artefacts from arithmetic $\lambda$-mixing between hydrophobic amino
acids and the non-hydrophobic phytate beads.  The Coulomb interaction
is inherited from the default via the \texttt{append} policy.

Intramolecular nonbonded pairs separated by one bond are excluded
(\texttt{excluded\_neighbours:~1}).  All intramolecular pairs within
each phytate molecule are excluded.
The aqueous medium is modelled with the temperature-dependent
dielectric permittivity of water at 298.15\,K.

\section{Monte Carlo moves}

Sampling is performed with the Metropolis criterion using
Faunus~\cite{Stenqvist2013}.  The following moves are applied per
sweep:

\begin{itemize}
\item \textbf{Rigid-body translation} of peptide chains in $xy$
      ($\delta_\text{max} = 10.0$\,\AA) and phytate molecules along $z$
      ($\delta_\text{max} = 5.0$\,\AA, 20 repeats).
\item \textbf{Pivot rotation} of peptide chains
      ($\delta\theta_\text{max} = 0.5$\,rad).
\item \textbf{Rigid-body rotation} of phytate molecules
      ($\delta\theta_\text{max} = 1.0$\,rad).
\item \textbf{Speciation (molecular swap)} for the 12 phytate
      titration reactions and 4 peptide atom-type swaps
      (10 repeats per draw).
\end{itemize}

\section{Potential of mean force}

The potential of mean force (PMF) for phytate as a function of
distance from the brush surface is computed using multi-walker
umbrella sampling with hard-wall windows.  The collective variable
is the $z$-position of the INO bead, spanning the half slit from the
grafted wall to the box centre, $[-L_z/2, 0]$, with 40 windows
6\,\AA\ wide and $\approx 2.4$\,\AA\ apart (strongly overlapping).

Each window starts from a common equilibrated state.  A harmonic
drive phase ($k_\text{drive} = 5.0$\,kJ\,mol$^{-1}$\,\AA$^{-2}$)
steers the phytate into the target window; the bias is then removed
for unbiased production sampling.  Per-window states are saved for
restart.

Adjacent-window histograms are stitched via the overlap-ratio method:
\[
\Delta F_{i \to i+1} = -k_BT \ln \frac{f_i}{f_{i+1}}
\]
where $f_i$ is the fraction of samples in window $i$ that fall within
the overlap region with window $i{+}1$.

PMFs are computed at pH~4, 6, 7, and 10 and ionic strengths
$I = 0.005$, $0.03$, and $0.15$\,M NaCl (5--150\,mM) to probe the
interplay of electrostatic screening and charge regulation.

\section{Free phytate diffusion}

The rotational and translational diffusion of a single phytate molecule in
bulk solution are estimated for each protonation state from a hydrodynamic
bead model, independently of the Monte Carlo simulations.  For a given net
charge, 20 conformers are generated with RDKit (ETKDGv3 embedding followed by
MMFF94 optimisation, with a UFF fallback where MMFF lacks parameters for the
phosphate anions)~\cite{RDKit}.  Each conformer is mapped to an all-atom bead
model --- one bead per atom with Bondi van der Waals radii --- and the
$6\times6$ Rotne--Prager--Yamakawa grand mobility matrix is evaluated with
PyGRPY~\cite{Zuk2018,Waszkiewicz2023}.  The diffusion tensors follow as
$\mathbf{D} = k_BT\,\boldsymbol{\mu}$ (water viscosity $\eta = 0.89$\,mPa\,s at
298.15\,K).  We report the isotropic rotational diffusion coefficient
$D_\text{r}$, the rotational correlation time $\tau_\text{r} = 1/(6D_\text{r})$,
the translational coefficient $D_\text{t}$, and the Stokes radius
$a_\text{H} = k_BT/(6\pi\eta D_\text{t})$, each averaged over the conformer
ensemble.

Protonation enters this estimate through geometry alone: deprotonation removes
the peripheral hydrogen beads and reshapes the phosphate arms, altering the
bare hydrodynamic shape.  Stick boundary conditions are assumed and no
ionic-atmosphere friction is included, so the values are order-of-magnitude
estimates of the free-phytate tumbling time ($\tau_\text{r}\approx0.1$\,ns).

\section{Position-dependent rotational mobility}

A relative, position-resolved rotational mobility of the phytate is extracted
from the Monte Carlo trajectories themselves.  Each accepted rigid-body
rotation contributes its squared angular step, and the mean-square angular
displacement per attempted move, $\langle\Delta\theta^2\rangle$, is accumulated
for every protonation state and pooled.  The ratio of this quantity to its bulk
value gives a rough rotational-diffusion profile,
\[
\frac{D(r)}{D_0} \approx
\frac{\langle\Delta\theta^2\rangle(r)}{\langle\Delta\theta^2\rangle_\text{bulk}},
\]
where the bulk reference is the mean over the windows nearest the box centre.
Monte Carlo steps carry no physical time scale, so only the relative profile
$D(r)/D_0$ is meaningful; it measures how strongly the brush hinders phytate
tumbling on approach.

\section{Zeta potential}

The electric-potential analysis records the laterally averaged charge density
$\sigma(z)$ in each slab alongside the mean electrostatic potential, the
screened sum
\[
\varphi(z) = \frac{2\pi l_B}{\kappa}\sum_{z'}\sigma(z')\,e^{-\kappa|z-z'|},
\]
with $l_B$ the Bjerrum length and $1/\kappa$ the Debye length of the medium, so
that $\varphi\to0$ in the bulk.  We identify the hydrodynamic shear plane with
the outer edge of the fixed brush charge --- the largest $r$ at which the
peptide charge density $|\sigma(z)|$ still exceeds 2\% of its peak --- and
estimate the $\zeta$-potential as $\varphi$ evaluated there.  The sign and
pH-dependence of this estimate are robust, whereas its magnitude depends on the
choice of shear plane ($\varphi$ is several-fold larger read inside the brush)
and on the planar, single-peptide geometry.

\bibliographystyle{unsrt}
\bibliography{references}
\end{document}
